

% Chapter 1 Introduction Introduce your project, putting the work in its academic or
% business context. Provide an overview of the work you’ve done.


\subsection{Document Structure}
\label{sec:document_structure}
The following document will discuss:
\begin{itemize}
    \item The background of the project
    \item The purpose of the project
    \item The scope of the project
    \item The aims and objectives of the project
    \item The requirements of the project
    \item The technologies, tools and datasets used in the project
    \item The evaluation methodology used in the project
\end{itemize}

\subsection{Project Purpose}
\label{sec:project_purpose}

The purpose of this project is to investigate the feasibility of using Natural Language Processing (NLP) technologies in addition to Large Language Models to improve on existing reccomendation systems. 
% Exploring using new information retrieval techniques to improve the quality of recommendations made to users.

This will be through the use of 'chatbots' which are able to understand and respond to natural language queries. 
The usage of chatbots is to provide a more natural and intuitive way for users to interact with the recommendation system.
It is hoped that this will lead to more accurate and relevant recommendations being made to users, 
% while also providing a more engaging and enjoyable user experience.
while also making the user feel heard and understood.

\subsection{Project Scope}
\label{sec:project_scope}
The project will focus on the following areas:
\begin{itemize}
    % \item Investigating the current state of the art in recommendation systems
    % \item Investigating the current state of the art in NLP and Large Language Models
    \item Implementing a chatbot that can interact with users and build a user profile (e.g. user document matrix)
    \item Integrating the chatbot with a recommendation system
    \item Implementing large language models to generate natural, free flowing conversations with the user
    \item Investigate guiding the user to explore and understand their preferences better
    \item Gathering explicit feedback from the user to improve recommendations
    \item Gathering implicit feedback from the user-item matrix to improve recommendations
    \item Working with the user to refine their preferences for a 'session' to pick a singular film
    \item Evaluating the performance of the chatbot in terms of reducing the cold start problem
    \item Evaluating the performance of the chatbot in terms of recommendation quality and user satisfaction
\end{itemize}

\subsection{Background and project topic discussion}
\label{sec:background}
Currently, recommendation systems typically use collaborative filtering or content-based filtering to make recommendations to users.
These systems are limited in their ability to understand the user's preferences and provide relevant recommendations.
Often, users are required to provide explicit feedback on items they like or dislike in order to receive accurate recommendations, and even then, the recommendations may not be relevant or useful to the user.

The use of chatbots in recommendation systems is a relatively new area of research, but it has the potential to revolutionize the way recommendations are made to users. By using chatbots, recommendation systems can provide a more natural and intuitive way for users to interact with the system, leading to more accurate and relevant recommendations being made.

Recent advancements in conversational AI have shown significant potential in enhancing recommendation systems. 
According to a survey on conversational recommender systems, integrating natural language processing and dialogue management techniques can improve user satisfaction and engagement. 
These systems can understand user preferences through natural conversations, making the recommendation process more dynamic and personalized.

I found a project in this area an interesting proposition, as it seeks to address a common problem that I, and many others face when trying to find something to watch on Netflix or Amazon Prime. 
While this project may not be viable for a commercial product due to its cost, I believe a proof of concept is viable and could show the potential of this technology. 
I seek to investigate and explore the possibilities of using large language models in conversational recommendation systems, and to evaluate the system compared to previous recommendation systems.

\subsection{Project Definition}
\label{sec:project_definition}
The project will aim to address the following research questions:
\begin{itemize}
    \item Can chatbots be used to improve the quality of recommendations made to users?
    \item Can natural, free flowing conversations that LLMs can have with users be integrated into information retrieval systems?
    \item Can chatbots be used to guide users to explore and understand their preferences better?
    \item Can chatbots be used to refine user preferences for a 'session' to pick a singular film?
    \item Can chatbots be used to reduce the cold start problem in recommendation systems?
    \item Can chatbots be used to improve user satisfaction with recommendation systems?
\end{itemize}

I chose a project that is a proof of concept, as I believe it has the potential to show the benefits of using chatbots in recommendation systems.
It involves both research and implementation, as I will need to investigate the current state of the art in recommendation systems and NLP, and then implement a system that integrates these technologies.




